\begin{theorem}
Let $1<p<\infty$.  Then there are maps $f\colon\mathbb{C}\to\mathbb{C}$ and
$v\colon\mathbb{C}\to\mathbb{R}$ with the following properties, where
$H=H_p$ is the extremal function (54) of \noderef{ExtremalH} and
$d=d(p)$ is the exponent of \noderef{DExp}:
\begin{enumerate}
\item[(a)] $f$ is continuous on $\mathbb{C}$, $f(H(\xi))=\xi$ for every
  $\xi\in\mathbb{C}$ and $H(f(z))=z$ for every $z\in\mathbb{C}$;
\item[(b)] $f$ is $C^\infty$ on $\mathbb{C}-\{0\}$;
\item[(c)] $f(tz)=t^{d}f(z)$ for every $t\ge0$ and every $z$, and
  $D^\nu f(tz)=t^{\,d-|\nu|}D^\nu f(z)$ for every $t>0$, every $z\neq0$ and
  every iterated partial derivative $D^\nu$ of \noderef{PartialWord};
\item[(d)] $v$ is continuous on $\mathbb{C}$, is $C^\infty$ on
  $\mathbb{C}-\{0\}$, and is $p$-harmonic on $\mathbb{C}$ in the sense of
  \noderef{PHarmonicOn};
\item[(e)] $\partial v/\partial z=f$ at every $z\neq0$, with
  $\partial/\partial z$ as in \noderef{DZ}.
\end{enumerate}

This is (55) of the paper together with the assertion of paper.tex:592 that it
defines a $p$-harmonic function of complex gradient $f$; the paper's note lists
that assertion among the facts which must be proved.  Every regularity
statement above is a conclusion and none is assumed: the paper's arrow runs
from the equation to the regularity.
\end{theorem}
\begin{proof}
Write $\mu=1/d$, $\varepsilon=\epsilon_2(1,p)$ for the coefficient of
\noderef{EpsCoef} appearing in \noderef{ExtremalH}, $\lambda_2=\lambda_2(1,p)$
for the exponent of \noderef{LambdaCoef}, and
$c=\tfrac12-\tfrac1p$, so that (18) at $n=1$ reads
\begin{equation}\tag{18}
  H_{\bar\xi}(\xi)=c\Bigl[\frac{\xi}{\bar\xi}H_\xi(\xi)
    +\frac{\bar\xi}{\xi}\overline{H_\xi(\xi)}\Bigr].
\end{equation}

\emph{Step 1: $d>0$ and $\lambda_2=\mu$.}  By \noderef{GammaExp},
$\gamma_1=\lambda_2(1,p)$.  By \noderef{DExpEqGammaOneInv}, $\gamma_1>0$ and
$\gamma_1\,d=1$.  Hence $\lambda_2>0$, $d=1/\lambda_2>0$ and
$\lambda_2=1/d=\mu$.

\emph{Step 2: the inverse, and its zero.}  By \noderef{ExtremalInjective}, $H$
is continuous and bijective and there is a continuous $f$ with
$f(H(\xi))=\xi$ and $H(f(z))=z$ for all $\xi,z\in\mathbb{C}$.  Fix this $f$;
this is (a).

The second clause of \noderef{ExtremalPolarForm} gives
$t\,H(\xi)=H(t^{d}\xi)$ for $t>0$; at $\xi=0$ and $t=2$ this reads
$2H(0)=H(0)$, so
\begin{equation}\tag{1}
  H(0)=0,\qquad\text{and hence}\qquad
  f(z)=0\iff z=0 .
\end{equation}
Indeed $f(0)=f(H(0))=0$, and if $f(z)=0$ then $z=H(f(z))=H(0)=0$.

\emph{Step 3: $H$ is the $k=2$ summand of (19) at $n=1$.}  Let
$A\colon\mathbb{Z}\to\mathbb{C}$ be $A(2)=1$ and $A(j)=0$ for $j\neq2$.  By
\noderef{Thm2Term} at $n=1$, $k=2$,
\[
  T(\xi):=\bigl(A(2)\xi^{2}+\varepsilon\,\overline{A(2)}\,\bar\xi^{\,2}\bigr)
    |\xi|^{\lambda_2+1-2}\xi^{-1}
  =\bigl(\xi^{2}+\varepsilon\bar\xi^{\,2}\bigr)|\xi|^{\lambda_2-1}\xi^{-1},
\]
$\varepsilon$ being real.  If $\xi\neq0$ then $\bar\xi=|\xi|^{2}/\xi$, so
$\bar\xi^{\,2}=|\xi|^{4}/\xi^{2}$ and
\[
  T(\xi)=\Bigl(\xi^{2}+\varepsilon\frac{|\xi|^{4}}{\xi^{2}}\Bigr)
    \frac{|\xi|^{\lambda_2-1}}{\xi}
  =\Bigl(\xi+\varepsilon\frac{|\xi|^{4}}{\xi^{3}}\Bigr)|\xi|^{\lambda_2-1}
  =\Bigl(\frac{\xi}{|\xi|}+\varepsilon\frac{|\xi|^{3}}{\xi^{3}}\Bigr)
    |\xi|^{\lambda_2}
  =H(\xi),
\]
by \noderef{ExtremalH} and $\lambda_2=\mu=1/d$ from Step 1.  At $\xi=0$ both
sides vanish: $T(0)=0$ because its first factor does, and $H(0)=0$ by (1).
Hence $H=T$ as functions on $\mathbb{C}$.

By \noderef{Thm2TermSolvesEq18} at $n=1$, $k=2$, whose hypotheses $1\le1$,
$1\le2$ and $1<p$ hold, $T$ is $C^\infty$ at every $\xi\neq0$ and solves (18)
at $n=1$ on $\{|\xi|<1\}$ in the sense of \noderef{SolvesEq18}.  Therefore
\begin{equation}\tag{2}
  H\ \text{is}\ C^\infty\ \text{on}\ \mathbb{C}-\{0\},
\end{equation}
the set $\mathbb{C}-\{0\}$ being open, and $H$ solves (18) on $\{|\xi|<1\}$.

\emph{Step 3$'$: weak and classical Wirtinger derivatives agree.}  We record
once, in general form, the passage between the two readings of the Wirtinger
operators, and use it twice below.  \emph{Let $\Omega\subseteq\mathbb{C}$ be
open, let $u\colon\mathbb{C}\to\mathbb{C}$ be differentiable at almost every
point of $\Omega$, and let $a,b$ be weak Wirtinger derivatives of $u$ on
$\Omega$ in the sense of \noderef{IsWeakWirtingerOn}.  Then $a=u_z$ and
$b=u_{\bar z}$ almost everywhere on $\Omega$, with $u_z,u_{\bar z}$ the
pointwise operators of \noderef{DZ} and \noderef{DZbar}.}  Indeed, by
\noderef{IsWeakWirtingerOn} there are $u_x,u_y$ with
$u,u_x,u_y\in L^1_{\mathrm{loc}}(\Omega)$ in the sense of \noderef{MemLpLocOn},
with $u_x$ and $u_y$ the distributional $\partial_0$- and
$\partial_1$-derivatives of $u$ on $\Omega$ in the sense of
\noderef{IsWeakPartialOn}, and with $a=\tfrac12(u_x-iu_y)$,
$b=\tfrac12(u_x+iu_y)$ at every point.  The first two are hypotheses (H1) and
(H2) of \noderef{WeakPartialEqFDerivAE} for $\Omega$, $u$ and the pair
$D_0=u_x$, $D_1=u_y$, and the almost-everywhere differentiability is its (H3);
so $u_x=\partial_0u$ and $u_y=\partial_1u$ almost everywhere on $\Omega$, where
$\partial_j=D^{[j]}$ is the classical partial derivative of
\noderef{PartialWord}, that is $\partial_0u=\mathrm{d}u[1]$ and
$\partial_1u=\mathrm{d}u[i]$; and $u_z,u_{\bar z}$ are by \noderef{DZ} and
\noderef{DZbar} the same two combinations
$\tfrac12(\mathrm{d}u[1]\mp i\,\mathrm{d}u[i])$.  Call this fact (W).  We note
for later use the two identities it rests on, valid at every point by
\noderef{DZ} and \noderef{DZbar}:
\begin{equation}\tag{2$'$}
  \partial_0u=u_z+u_{\bar z},\qquad \partial_1u=i\,(u_z-u_{\bar z}).
\end{equation}

\emph{Step 4: (18) at every point of the punctured unit disc.}  By (2) the map
$\xi\mapsto\mathrm{d}H(\xi)$ is continuous on $\mathbb{C}-\{0\}$, so $H_\xi$ and
$H_{\bar\xi}$ of \noderef{DZ} and \noderef{DZbar}, being the fixed combinations
$\tfrac12(\mathrm{d}H[1]\mp i\,\mathrm{d}H[i])$, are continuous there, and so
are $\xi\mapsto\xi/\bar\xi$ and $\xi\mapsto\bar\xi/\xi$.  Hence both sides of
(18) are continuous on the open set $\{0<|\xi|<1\}$.  By Step 3 and
\noderef{SolvesEq18} there are weak Wirtinger derivatives $a_H,b_H$ of $H$ on
$\mathbb{B}=\{|\xi|<1\}$ and a null set $Z_0$ such that the identity (18) holds
for the pair $a_H,b_H$ at every $\xi\in\mathbb{B}-Z_0$ with $\xi\neq0$; and by
(W), applied to $u=H$ on the open set $\mathbb{B}$ --- legitimate because by
(2) $H$ is differentiable at every point of $\mathbb{B}$ except possibly $0$,
a null set --- there is a null set $Z_1$ with $a_H=H_\xi$ and
$b_H=H_{\bar\xi}$ on $\mathbb{B}-Z_1$.  So, with $Z=Z_0\cup Z_1$, a null set,
(18) holds for $H_\xi,H_{\bar\xi}$ at every $\xi\notin Z$ with $|\xi|<1$ and
$\xi\neq0$.  A null set has empty interior, so
$\{0<|\xi|<1\}-Z$ is dense in $\{0<|\xi|<1\}$, and two continuous functions on
an open set agreeing on a dense subset agree everywhere on it.  Hence (18)
holds at every $\xi$ with $0<|\xi|<1$.

\emph{Step 5: (18) at every $\xi\neq0$.}  Let $\lambda>0$ and put
$t=\lambda^{\mu}>0$, so that $t^{d}=\lambda^{\mu d}=\lambda$; the second clause
of \noderef{ExtremalPolarForm} then gives
\begin{equation}\tag{3}
  H(\lambda\xi)=\lambda^{\mu}H(\xi)\qquad(\xi\in\mathbb{C},\ \lambda>0).
\end{equation}
Let $m(\xi)=\lambda\xi$.  Then $\mathrm{d}m(\xi)[1]=\lambda$ and
$\mathrm{d}m(\xi)[i]=\lambda i$, so by \noderef{DZ} and \noderef{DZbar},
\[
  \partial_\xi m=\tfrac12(\lambda-i\cdot\lambda i)=\lambda,\qquad
  \partial_{\bar\xi}m=\tfrac12(\lambda+i\cdot\lambda i)=0 .
\]
Fix $\xi_0\neq0$ and choose $\lambda$ with $0<\lambda<1/|\xi_0|$, so that
$0<|\lambda\xi_0|<1$.  By (2), $H$ is differentiable at $\lambda\xi_0=m(\xi_0)$,
and $m$ is differentiable at $\xi_0$, so \noderef{WirtingerChainRule} applies
to the pair $(H,m)$ at $\xi_0$ and gives
\[
  \partial_\xi(H\circ m)(\xi_0)=\lambda\,H_\xi(\lambda\xi_0),\qquad
  \partial_{\bar\xi}(H\circ m)(\xi_0)=\lambda\,H_{\bar\xi}(\lambda\xi_0),
\]
using $\partial_\xi m=\lambda$, $\partial_{\bar\xi}m=0$ and
$\overline{\lambda}=\lambda$.  On the other hand $H\circ m=\lambda^{\mu}H$ as
functions by (3), and $\partial_\xi$, $\partial_{\bar\xi}$ are
$\mathbb{R}$-linear in the function, so
$\partial_\xi(H\circ m)(\xi_0)=\lambda^{\mu}H_\xi(\xi_0)$ and likewise for
$\partial_{\bar\xi}$.  Comparing, and dividing by $\lambda\neq0$,
\begin{equation}\tag{4}
  H_\xi(\lambda\xi_0)=\lambda^{\mu-1}H_\xi(\xi_0),\qquad
  H_{\bar\xi}(\lambda\xi_0)=\lambda^{\mu-1}H_{\bar\xi}(\xi_0).
\end{equation}
Since $\lambda>0$ is real, $\lambda\xi_0/\overline{\lambda\xi_0}=\xi_0/\bar\xi_0$
and $\overline{\lambda\xi_0}/(\lambda\xi_0)=\bar\xi_0/\xi_0$.  Substituting (4)
into (18) at the point $\lambda\xi_0$, which is legitimate by Step 4, and
dividing by $\lambda^{\mu-1}\neq0$, we obtain (18) at $\xi_0$.  Hence
\begin{equation}\tag{5}
  \text{(18) holds at every }\xi\neq0,
\end{equation}
where (18) is read for the pointwise derivatives $H_\xi$, $H_{\bar\xi}$.

\emph{Step 5$'$: $H$ solves (18) on $\mathbb{C}$ in the sense of
\noderef{SolvesEq18}.}  We must produce weak Wirtinger derivatives of $H$ on
$\mathbb{C}$ and check (18) for them almost everywhere.

First a growth bound near the origin.  Let $0<|\xi|<1$ and put
$\xi_0=\xi/|\xi|$ and $\lambda=|\xi|$, so that $|\xi_0|=1$, $\lambda>0$,
$\lambda<1=1/|\xi_0|$ and $\lambda\xi_0=\xi$; these are exactly the conditions
under which (4) was established, so
$H_\xi(\xi)=\lambda^{\mu-1}H_\xi(\xi_0)$ and
$H_{\bar\xi}(\xi)=\lambda^{\mu-1}H_{\bar\xi}(\xi_0)$.  By (2) the functions
$H_\xi$ and $H_{\bar\xi}$ are continuous on $\mathbb{C}-\{0\}$, hence bounded
on the compact unit circle, say by $C_0$.  Therefore
\[
  |H_\xi(\xi)|\le C_0|\xi|^{\mu-1},\qquad
  |H_{\bar\xi}(\xi)|\le C_0|\xi|^{\mu-1}\qquad(0<|\xi|<1),
\]
whence by (2$'$), for $j=0,1$,
\begin{equation}\tag{5$'$}
  |\partial_jH(\xi)|\le 2C_0\,|\xi|^{\mu-1}\qquad(0<|\xi|<1).
\end{equation}

Next, $H$ and $\partial_0H,\partial_1H$ lie in
$L^1_{\mathrm{loc}}(\mathbb{C})$ in the sense of \noderef{MemLpLocOn}.  For $H$
this holds because $H$ is continuous on $\mathbb{C}$ by Step 2, hence bounded
on every compact set, which has finite measure.  For $\partial_jH$:  it is
continuous on $\mathbb{C}-\{0\}$ by (2), hence measurable off the null set
$\{0\}$ and so measurable.  Let $\kappa\subseteq\mathbb{C}$ be compact and
split $\kappa=\kappa_1\cup\kappa_2$ with
$\kappa_1=\kappa\cap\{|\xi|\le\tfrac12\}$ and
$\kappa_2=\kappa\cap\{|\xi|\ge\tfrac12\}$.  On $\kappa_1$, by (5$'$) and polar
coordinates,
\[
  \int_{\kappa_1}|\partial_jH|\,d\sigma
  \le 2C_0\int_{\{|\xi|<1/2\}}|\xi|^{\mu-1}d\sigma(\xi)
  =4\pi C_0\int_0^{1/2}r^{\mu-1}\,r\,dr
  =4\pi C_0\int_0^{1/2}r^{\mu}\,dr<\infty,
\]
because $\mu>0>-1$ by Step 1.  The set $\kappa_2$ is closed and bounded, hence
compact, and contained in $\mathbb{C}-\{0\}$, where $\partial_jH$ is
continuous; so $\partial_jH$ is bounded on $\kappa_2$, which has finite
measure, and $\int_{\kappa_2}|\partial_jH|\,d\sigma<\infty$.  Adding the two
bounds gives $\int_\kappa|\partial_jH|\,d\sigma<\infty$.

Now fix $j\in\{0,1\}$ and apply \noderef{WeakPartialAcrossPoint} with
$\Omega=\mathbb{C}$, $S=\{0\}$, this $j$, $g=H$, $v=\partial_jH$ and
$\beta=0$.  Its hypotheses hold: $\mathbb{C}$ is open and $S\subseteq\mathbb{C}$
with $S\cap\kappa$ finite for every compact $\kappa$, $S$ being a singleton;
$0\le\beta=0<1$; $g$ is $C^1$ on $\mathbb{C}-S$ by (2), a function $C^\infty$
on an open set being $C^1$ there; $D^{[j]}g=v$ on
$\mathbb{C}-S$ by the definition of $v$; $g,v\in L^1_{\mathrm{loc}}(\mathbb{C})$
by the previous paragraph; and the growth hypothesis holds at the only point
$0$ of $S$ with $\rho=1$ and the constant
$M=\sup_{|\xi|\le1}|H(\xi)|<\infty$, finite because $H$ is continuous on the
compact set $\{|\xi|\le1\}$, since $\|H(\xi)\|\le M=M|\xi-0|^{-0}$ for
$0<|\xi|<1$.  Therefore $\partial_jH$ is the distributional
$\partial_j$-derivative of $H$ on $\mathbb{C}$ in the sense of
\noderef{IsWeakPartialOn}.

Consequently the three clauses of \noderef{IsWeakWirtingerOn} hold for $H$ on
$\mathbb{C}$ with $u_x=\partial_0H$, $u_y=\partial_1H$ and the pair
$a=H_\xi$, $b=H_{\bar\xi}$: clause (a) is the local integrability just proved,
clause (b) is the last paragraph, and clause (c) is (2$'$) solved for $H_\xi$
and $H_{\bar\xi}$, that is
$H_\xi=\tfrac12(\partial_0H-i\partial_1H)$ and
$H_{\bar\xi}=\tfrac12(\partial_0H+i\partial_1H)$, which hold at every point by
\noderef{DZ} and \noderef{DZbar}.  So $H_\xi$ and $H_{\bar\xi}$ are themselves
weak Wirtinger derivatives of $H$ on $\mathbb{C}$, and by (5) the identity
(18) holds for them at every $\xi\neq0$, hence at almost every
$\xi\in\mathbb{C}$ with $\xi\neq0$.  This is exactly the assertion that $H$
solves (18) at $n=1$ on $\mathbb{C}$ in the sense of \noderef{SolvesEq18};
record it as
\begin{equation}\tag{5$''$}
  H\ \text{solves (18) at }n=1\ \text{on }\mathbb{C}\ \text{in the sense of
  \noderef{SolvesEq18}} .
\end{equation}

\emph{Step 6: $H_\xi\neq0$ off the origin.}  Let $\xi\neq0$.  By
\noderef{ExtremalH} and $|\xi|>0$,
\[
  H(\xi)=\xi\,|\xi|^{\mu-1}+\varepsilon\,|\xi|^{\mu+3}\xi^{-3}
  =:T_1(\xi)+T_2(\xi).
\]
By \noderef{DZ} applied to the identity map, $\partial_\xi\xi=1$ and
$\partial_{\bar\xi}\xi=0$.  By \noderef{WirtingerModulusPowerRule},
$\partial_\xi(|\xi|^{m})=\tfrac m2|\xi|^{m}/\xi$ for every real $m$, and these
functions are $C^\infty$ at $\xi$.  By \noderef{WirtingerPowerRule} applied to
the identity map on the open set $\mathbb{C}-\{0\}$, where it is $C^\infty$ and
non-vanishing, with the integer exponent $-3$,
$\partial_\xi(\xi^{-3})=-3\xi^{-4}$.  Hence, by
\noderef{WirtingerProductRule},
\[
  \partial_\xi T_1=|\xi|^{\mu-1}+\xi\cdot\frac{\mu-1}{2}
    \frac{|\xi|^{\mu-1}}{\xi}=\frac{\mu+1}{2}|\xi|^{\mu-1},
\]
\[
  \partial_\xi T_2
  =\varepsilon\Bigl[\frac{\mu+3}{2}\frac{|\xi|^{\mu+3}}{\xi}\xi^{-3}
    -3|\xi|^{\mu+3}\xi^{-4}\Bigr]
  =\varepsilon\frac{\mu-3}{2}|\xi|^{\mu+3}\xi^{-4}
  =\varepsilon\frac{\mu-3}{2}|\xi|^{\mu-1}\frac{|\xi|^{4}}{\xi^{4}} ,
\]
where the last step uses $|\xi|^{\mu+3}=|\xi|^{\mu-1}|\xi|^{4}$.  Since
$\bigl||\xi|^{4}/\xi^{4}\bigr|=1$,
\begin{equation}\tag{6}
  |H_\xi(\xi)|\ \ge\ |\xi|^{\mu-1}\cdot
    \frac{(\mu+1)-|\varepsilon|\,|\mu-3|}{2} .
\end{equation}
We check that the last numerator is positive.  By
\noderef{ExtremalEpsBound}, $\varepsilon=(1-d)/(1+3d)$, and $d>0$ by Step 1, so
$1+3d>0$ and $|\varepsilon|=|1-d|/(1+3d)$.  With $\mu=1/d$,
\[
  \mu+1=\frac{1+d}{d},\qquad |\mu-3|=\frac{|1-3d|}{d},
\]
so $(\mu+1)-|\varepsilon||\mu-3|>0$ is equivalent, after multiplying by the
positive number $d(1+3d)$, to
\[
  (1+d)(1+3d)>|1-d|\,|1-3d|=|(1-d)(1-3d)| ,
\]
that is, to $1+4d+3d^{2}>|1-4d+3d^{2}|$.  This holds because
$(1+4d+3d^{2})-(1-4d+3d^{2})=8d>0$ and
$(1+4d+3d^{2})+(1-4d+3d^{2})=2+6d^{2}>0$.  Hence the numerator in (6) is
positive and
\begin{equation}\tag{7}
  H_\xi(\xi)\neq0\qquad(\xi\neq0).
\end{equation}

\emph{Step 7: $\det\mathrm{d}H=J_H$, and $J_H>0$ off the origin.}  Let
$F\colon\mathbb{C}\to\mathbb{C}$ be differentiable at $w$ and put
$L=\mathrm{d}F(w)$, $a=F_z(w)$, $b=F_{\bar z}(w)$.  By \noderef{DZ} and
\noderef{DZbar}, $a=\tfrac12(L[1]-iL[i])$ and $b=\tfrac12(L[1]+iL[i])$, so
$L[1]=a+b$ and $L[i]=i(a-b)$.  In the basis $\{1,i\}$ the matrix of $L$ has
columns $(\operatorname{Re}(a+b),\operatorname{Im}(a+b))$ and
$(-\operatorname{Im}(a-b),\operatorname{Re}(a-b))$, so
\[
  \det L=\operatorname{Re}(a+b)\operatorname{Re}(a-b)
    +\operatorname{Im}(a+b)\operatorname{Im}(a-b)
  =\operatorname{Re}\bigl((a+b)\overline{(a-b)}\bigr)=|a|^{2}-|b|^{2},
\]
because $b\bar a-a\bar b=b\bar a-\overline{b\bar a}$ is purely imaginary.  By
\noderef{Jac} this says
\begin{equation}\tag{8}
  \det\mathrm{d}F(w)=J_F(w)=|F_z(w)|^{2}-|F_{\bar z}(w)|^{2}.
\end{equation}

Now let $\xi\neq0$.  Taking moduli in (5) and using
$|\xi/\bar\xi|=|\bar\xi/\xi|=1$ and
$|\overline{H_\xi}|=|H_\xi|$,
\[
  |H_{\bar\xi}(\xi)|\le2|c|\,|H_\xi(\xi)|=\frac{|p-2|}{p}|H_\xi(\xi)| .
\]
For $p>1$ we have $|p-2|<p$, since $p-2<p$ and $-(p-2)=2-p<p$ is
$1<p$; so $|p-2|/p<1$ and, by (7) and \noderef{Jac},
\begin{equation}\tag{9}
  J_H(\xi)=|H_\xi(\xi)|^{2}-|H_{\bar\xi}(\xi)|^{2}
  \ \ge\ \Bigl(1-\frac{(p-2)^{2}}{p^{2}}\Bigr)|H_\xi(\xi)|^{2}>0
  \qquad(\xi\neq0).
\end{equation}

\emph{Step 8: $f$ is $C^\infty$ off the origin, with $J_f\neq0$ there.}  Fix
$z_0\neq0$ and put $\xi_0=f(z_0)$.  By (1), $\xi_0\neq0$, and
$H(\xi_0)=H(f(z_0))=z_0$ by Step 2.  By (2), $H$ is $C^\infty$ on a
neighbourhood of $\xi_0$, and by (8) and (9) its differential
$\mathrm{d}H(\xi_0)$ has non-zero determinant, hence is invertible.  By the
inverse function theorem there are open sets $N\ni\xi_0$ and $M\ni z_0$ such
that $H$ maps $N$ bijectively onto $M$ and the inverse
$h\colon M\to N$ of this restriction is $C^\infty$.

Since $f$ is continuous and $f(z_0)=\xi_0\in N$, there is an open $M'$ with
$z_0\in M'\subseteq M$ and $f(M')\subseteq N$.  For $z\in M'$ we have
$H(f(z))=z$ with $f(z)\in N$ and $z\in M$, so $f(z)=h(H(f(z)))=h(z)$.  Thus
$f=h$ on $M'$, and $f$ is $C^\infty$ on a neighbourhood of $z_0$.  As
$z_0\neq0$ was arbitrary, this proves (b).

Moreover, differentiating $H\circ f=\mathrm{id}$ on $M'$ at $z_0$ gives
$\mathrm{d}H(\xi_0)\circ\mathrm{d}f(z_0)=\mathrm{id}$, so $\mathrm{d}f(z_0)$ is
invertible and, by (8) applied to $f$,
\begin{equation}\tag{10}
  J_f(z)=\det\mathrm{d}f(z)\neq0\qquad(z\neq0).
\end{equation}

\emph{Step 9: homogeneity.}  By (b), $f$ is $C^\infty$ on
$\{0\}^{c}=\mathbb{C}-\{0\}$, so \noderef{ExtremalHomogeneous} applies with the
two inverse relations of Step 2, and yields exactly the two clauses of (c).
Nothing in its hypotheses refers to equation (12), so this step may be taken
before Step 11.

\emph{Step 10: $f\in W^{1,1}_{\mathrm{loc}}(\mathbb{C})$.}  By the second
clause of (c) at $|\nu|=1$, that is for the one-letter words $w=[j]$,
\[
  \partial_jf(tz)=t^{\,d-1}\partial_jf(z)\qquad(t>0,\ z\neq0,\ j=0,1),
\]
$\partial_j=D^{[j]}$ being the classical partial derivative of
\noderef{PartialWord}.  By (b) the functions $\partial_0f,\partial_1f$ are
continuous on $\mathbb{C}-\{0\}$, hence bounded on the compact unit circle, say
by $C_1$; taking $t=|z|$ and the point $z/|z|$ of that circle,
\begin{equation}\tag{10$'$}
  |\partial_jf(z)|\le C_1|z|^{\,d-1}\qquad(z\neq0,\ j=0,1).
\end{equation}
Hence $\partial_jf\in L^1_{\mathrm{loc}}(\mathbb{C})$ in the sense of
\noderef{MemLpLocOn}: it is continuous off the null set $\{0\}$ and so
measurable, and for compact $\kappa\subseteq B(0,R)$, by (10$'$) and polar
coordinates,
\[
  \int_\kappa|\partial_jf|\,d\sigma
  \le C_1\int_{B(0,R)}|z|^{\,d-1}d\sigma(z)
  =2\pi C_1\int_0^{R}r^{\,d-1}r\,dr=2\pi C_1\int_0^{R}r^{\,d}\,dr<\infty,
\]
because $d>0>-1$ by Step 1.  Also $f\in L^1_{\mathrm{loc}}(\mathbb{C})$, being
continuous on $\mathbb{C}$ by (a) and hence bounded on each compact set, which
has finite measure.

Fix $j\in\{0,1\}$ and apply \noderef{WeakPartialAcrossPoint} with
$\Omega=\mathbb{C}$, $S=\{0\}$, this $j$, $g=f$, $v=\partial_jf$ and
$\beta=0$: $\mathbb{C}$ is open, $S$ is a singleton contained in it so
$S\cap\kappa$ is finite for every compact $\kappa$, $0\le\beta=0<1$, $g$ is
$C^1$ on $\mathbb{C}-S$ by (b), a function $C^\infty$ on an open set being
$C^1$ there, $D^{[j]}g=v$ there by the definition of $v$,
$g$ and $v$ are in $L^1_{\mathrm{loc}}(\mathbb{C})$ by the last paragraph, and
the growth hypothesis holds at $0$ with $\rho=1$ and
$M=\sup_{|z|\le1}|f(z)|<\infty$, finite by continuity of $f$ on the compact
disc $\{|z|\le1\}$.  So $\partial_jf$ is the distributional
$\partial_j$-derivative of $f$ on $\mathbb{C}$ in the sense of
\noderef{IsWeakPartialOn}.

Therefore the family $D$ with $D_{()}=f$, $D_{(j)}=\partial_jf$ for $j=0,1$ and
$D_w=0$ for $|w|\ge2$ witnesses the three clauses of \noderef{MemWLocOn} for
$k=1$, $s=1$ on $\mathbb{C}$: clause (a) is $D_{()}=f$; clause (b) is the local
integrability of $f$ and of $\partial_0f,\partial_1f$ just proved, the
functions $D_w$ with $|w|\ge2$ being irrelevant since $|w|\le k=1$; and clause
(c), which concerns only the empty word, is the previous paragraph.  Hence
\begin{equation}\tag{10$''$}
  f\in W^{1,1}_{\mathrm{loc}}(\mathbb{C})\quad\text{in the sense of
  \noderef{MemWLocOn}} .
\end{equation}

\emph{Step 11: $f$ solves (12) on $\mathbb{C}$.}  We apply
\noderef{InverseSolvesEq12} with $V=\mathbb{C}$ and the pair $(H,f)$.  Its
hypotheses hold: $\mathbb{C}$ is open; $H(\mathbb{C})=\mathbb{C}$ is open, $H$
being surjective by Step 2; $H$ is $C^1$ on $\mathbb{C}-\{0\}$ by (2); $H$
solves (18) at $n=1$ on $\mathbb{C}$ in the sense of \noderef{SolvesEq18} by
(5$''$); $f(H(\xi))=\xi$ for every $\xi$ by Step 2; by (b) and (10), $f$ is
differentiable with $J_f\neq0$ at every $z\neq0$, hence at almost every $z$,
the origin being a null set; and its hypothesis (H6),
$f\in W^{1,1}_{\mathrm{loc}}(H(\mathbb{C}))=W^{1,1}_{\mathrm{loc}}(\mathbb{C})$,
is (10$''$), the image $H(\mathbb{C})$ being $\mathbb{C}$ by Step 2.  The
conclusion is that $f$ solves (12) on $\mathbb{C}$ in the sense of
\noderef{SolvesEq12}: there are weak Wirtinger derivatives $a_f,b_f$ of $f$ on
$\mathbb{C}$ and a null set off which
\[
  b_f=\Bigl(\frac1p-\frac12\Bigr)
    \Bigl[\frac{\overline{f}}{f}a_f+\frac{f}{\overline{f}}
      \overline{a_f}\Bigr]
\]
at the points where $f\neq0$.  By (W) of Step 3$'$, applied to $u=f$ on the open set
$\Omega=\mathbb{C}$ --- legitimate because $f$ is differentiable at every point
of $\mathbb{C}-\{0\}$ by (b) and $\{0\}$ is null --- we have $a_f=f_z$ and
$b_f=f_{\bar z}$ almost everywhere, so, the union of two null sets being null,
\begin{equation}\tag{11}
  f_{\bar z}=\Bigl(\frac1p-\frac12\Bigr)
    \Bigl[\frac{\overline{f}}{f}f_z+\frac{f}{\overline{f}}
      \overline{f_z}\Bigr]
\end{equation}
almost everywhere where $f\neq0$, the derivatives now being the pointwise ones
of \noderef{DZ} and \noderef{DZbar}.

\emph{Step 12: (12) holds at every $z\neq0$, and $f_{\bar z}$ is real there.}  By
(1), $f$ does not vanish on the open set $\mathbb{C}-\{0\}$, and by (b) it is
$C^\infty$ there, so, as in Step 4, both sides of (11) are continuous on
$\mathbb{C}-\{0\}$ while (11) holds off a null set; the complement of a null set
is dense, so (11) holds at every $z\neq0$.  Writing
$a=(\overline{f}/f)f_z$ we have $(f/\overline{f})\overline{f_z}=\bar a$, so the
bracket in (11) is $a+\bar a$, which is real, and $1/p-1/2$ is real.  Hence
\begin{equation}\tag{12}
  \overline{f_{\bar z}(z)}=f_{\bar z}(z)\qquad(z\neq0).
\end{equation}

\emph{Step 13: Euler's identity.}  Let $F$ be differentiable at $w$ and
$L=\mathrm{d}F(w)$, $a=F_z(w)$, $b=F_{\bar z}(w)$.  By the computation in
Step 7, $L[1]=a+b$ and $L[i]=i(a-b)$, so for $h=h_1+ih_2$ with $h_1,h_2$ real,
\[
  L[h]=h_1(a+b)+h_2\,i(a-b)=a(h_1+ih_2)+b(h_1-ih_2)=a\,h+b\,\bar h .
\]
Fix $z\neq0$ and put $\psi(t)=f(tz)$ for $t>0$.  By (b) and the chain rule,
$\psi$ is differentiable and $\psi'(1)=\mathrm{d}f(z)[z]=z f_z(z)+\bar z
f_{\bar z}(z)$.  By the first clause of (c), $\psi(t)=t^{d}f(z)$ for $t>0$, so
$\psi'(1)=d\,f(z)$.  Hence
\begin{equation}\tag{13}
  z f_z(z)+\bar z f_{\bar z}(z)=d\,f(z)\qquad(z\neq0).
\end{equation}

\emph{Step 14: the potential.}  By Step 1, $d+1>0$.  Define
\[
  v(z)=\frac{2}{d+1}\operatorname{Re}\bigl(z f(z)\bigr)\qquad(z\in\mathbb{C}),
\]
a real-valued function, continuous on $\mathbb{C}$ because $f$ is.  Let
$\varphi_v$ denote $v$ regarded as complex-valued, so that
\[
  \varphi_v(z)=\frac{1}{d+1}\Bigl(z f(z)+\overline{z f(z)}\Bigr).
\]
Since $f$ is $C^\infty$ on $\mathbb{C}-\{0\}$ and conjugation is a real-linear
isometry, $\varphi_v$ is $C^\infty$ on $\mathbb{C}-\{0\}$.  Fix $z\neq0$.  By
$\partial_z z=1$, $\partial_{\bar z}z=0$ and
\noderef{WirtingerProductRule},
\[
  \partial_z\bigl(zf(z)\bigr)=f(z)+zf_z(z),\qquad
  \partial_{\bar z}\bigl(zf(z)\bigr)=z f_{\bar z}(z),
\]
and by \noderef{WirtingerConjRule} and (12),
\[
  \partial_z\overline{\bigl(zf(z)\bigr)}
  =\overline{\partial_{\bar z}\bigl(zf(z)\bigr)}
  =\overline{z f_{\bar z}(z)}=\bar z\,f_{\bar z}(z).
\]
Adding and using (13),
\begin{equation}\tag{14}
  \partial_z\varphi_v(z)=\frac{1}{d+1}
    \bigl(f(z)+zf_z(z)+\bar z f_{\bar z}(z)\bigr)
  =\frac{1}{d+1}\bigl(f(z)+d\,f(z)\bigr)=f(z)\qquad(z\neq0),
\end{equation}
which is (e).

\emph{Step 15: local integrability of the derivatives of $|f|^{p-2}f$.}  Put
$w=|f|^{p-2}f$ and $\kappa=(p-1)d>0$.  By (1) and (b), $f$ is $C^\infty$ and
non-vanishing on $\mathbb{C}-\{0\}$, and $\zeta\mapsto|\zeta|^{p-2}$ is
$C^\infty$ on $\mathbb{C}-\{0\}$ by \noderef{WirtingerModulusPowerRule}, so $w$
is $C^\infty$ on $\mathbb{C}-\{0\}$.  For $t>0$ and $z\neq0$, the first clause
of (c) gives $f(tz)=t^{d}f(z)$ with $t^{d}>0$, so
\[
  w(tz)=\bigl|t^{d}f(z)\bigr|^{p-2}t^{d}f(z)=t^{(p-2)d}t^{d}w(z)
  =t^{\kappa}w(z).
\]
Fix $t>0$ and $j\in\{0,1\}$.  Differentiating the identity
$w(tz)=t^{\kappa}w(z)$ in $z$ on the open set $\mathbb{C}-\{0\}$, the left side
has $j$-th partial $t\,(\partial_jw)(tz)$ by the chain rule and the right side
has $t^{\kappa}\partial_jw(z)$, so
\[
  (\partial_jw)(tz)=t^{\kappa-1}\partial_jw(z)\qquad(z\neq0,\ t>0).
\]
The function $\partial_jw$ is continuous on $\mathbb{C}-\{0\}$, hence bounded on
the unit circle, say by $C$.  For $z\neq0$, applying the last display with
$t=|z|$ and the point $z/|z|$ of the unit circle,
\[
  |\partial_jw(z)|=|z|^{\kappa-1}\bigl|\partial_jw(z/|z|)\bigr|
  \le C|z|^{\kappa-1}.
\]
Let $K\subseteq\mathbb{C}$ be compact and choose $R$ with
$K\subseteq B(0,R)$.  In polar coordinates
\[
  \int_{B(0,R)}|z|^{\kappa-1}\,\mathrm{d}\sigma(z)
  =2\pi\int_0^{R}r^{\kappa-1}\,r\,\mathrm{d}r
  =2\pi\int_0^{R}r^{\kappa}\,\mathrm{d}r<\infty,
\]
because $\kappa>0>-1$.  Hence $\partial_jw$, which is continuous off the null
set $\{0\}$ and so measurable, is integrable on $K$; that is,
$\partial_jw\in L^1_{\mathrm{loc}}(\mathbb{C})$ in the sense of
\noderef{MemLpLocOn}.

\emph{Step 16: $v$ is $p$-harmonic.}  We apply \noderef{Eq12ImpliesPHarmonic}
with $\Omega=\mathbb{C}$, $S=\{0\}$, and the pair $(f,v)$.  Its hypotheses hold:
$S\cap K\subseteq\{0\}$ is finite for every compact $K$; $f(z)=0$ exactly for
$z\in S$ by (1); $f$ is continuous by (a) and $C^\infty$ on
$\mathbb{C}-\{0\}$ by (b); $f$ solves (12) on $\mathbb{C}$ by Step 11; $v$ is
continuous on $\mathbb{C}$ and $C^\infty$ on $\mathbb{C}-\{0\}$ by Step 14;
$\partial v/\partial z=f$ on $\mathbb{C}-\{0\}$ by (14); $v$ and $f$ are
continuous, hence bounded on compact sets, hence in
$L^p_{\mathrm{loc}}(\mathbb{C})$ in the sense of \noderef{MemLpLocOn}; and the
two first-order partial derivatives of $|f|^{p-2}f$ are in
$L^1_{\mathrm{loc}}(\mathbb{C})$ by Step 15.  The conclusion is that $v$ is
$p$-harmonic on $\mathbb{C}$ in the sense of \noderef{PHarmonicOn}, which
together with the continuity of $v$ on $\mathbb{C}$ and its smoothness on
$\mathbb{C}-\{0\}$, both established in Step 14, is (d).

Steps 2, 8, 9, 14 and 16 establish (a), (b), (c), (e) and (d) respectively, so
the pair $(f,v)$ has all the stated properties.
\end{proof}
